\chapter{Manual de despliegue}\label{cap:manual_despliegue}

Este apéndice describe los pasos necesarios para desplegar el backend propio (API Spring Boot y base de datos PostgreSQL) en Railway~\cite{railway}, así como la configuración necesaria en el cliente React Native y en el proyecto de Firebase para que la aplicación completa funcione en un entorno nuevo.

\section{Requisitos previos}

\begin{itemize}
    \item Cuenta en Railway~\cite{railway} vinculada al repositorio de GitHub~\cite{github} del backend.
    \item Cuenta de Firebase con un proyecto creado, con los servicios \textbf{Storage} y \textbf{Cloud Messaging} habilitados (ya no es necesario habilitar Firestore ni Authentication).
    \item Clave de API de Google Maps Platform~\cite{googlemapsapi} con las APIs de \emph{Maps SDK} y \emph{Directions} activadas.
    \item Java 17 o superior y Node.js instalados en el entorno de desarrollo local.
\end{itemize}

\section{Despliegue del backend en Railway}

\begin{enumerate}
    \item Crear un nuevo proyecto en Railway y añadir dos servicios: uno de tipo \textbf{PostgreSQL} (Railway lo aprovisiona automáticamente) y otro de tipo \textbf{despliegue desde GitHub}, apuntando al repositorio del backend.
    \item Railway detecta automáticamente que se trata de un proyecto Spring Boot (Maven/Gradle) y genera la imagen de despliegue sin necesidad de un \texttt{Dockerfile} propio, aunque puede añadirse uno para tener un control más explícito del proceso de construcción.
    \item Configurar las siguientes variables de entorno en el servicio del backend:
    \begin{itemize}
        \item \texttt{SPRING\_DATASOURCE\_URL}, \texttt{SPRING\_DATASOURCE\_USERNAME} y \texttt{SPRING\_DATASOURCE\_PASSWORD}: Railway las expone automáticamente al enlazar el servicio de PostgreSQL con el del backend.
        \item \texttt{JWT\_SECRET}: clave secreta utilizada para firmar los tokens JWT (Sección~\ref{cap:1_iteracion}). Debe ser una cadena aleatoria de suficiente longitud y no debe compartirse ni subirse al repositorio.
        \item \texttt{JWT\_EXPIRATION}: tiempo de validez del token, en milisegundos.
        \item \texttt{FIREBASE\_CREDENTIALS}: contenido (o ruta) del fichero de credenciales de servicio de Firebase, necesario para que el backend pueda usar el Firebase Admin SDK~\cite{fcmadminsdk} y así disparar notificaciones \emph{push}.
    \end{itemize}
    \item Al desplegar, Spring Boot ejecuta automáticamente las migraciones del esquema de base de datos (Apéndice~\ref{cap:esquemaBD}) sobre la instancia de PostgreSQL antes de arrancar la API.
    \item Railway asigna un dominio público HTTPS al backend (por ejemplo, \texttt{https://\{proyecto\}.up.railway.app}), que será la URL base que consumirá la aplicación cliente.
    \item Cada nuevo \emph{commit} en la rama principal del repositorio dispara automáticamente un nuevo despliegue (integración continua), sin intervención manual.
\end{enumerate}

\section{Documentación de la API}

Una vez desplegado el backend, la documentación interactiva de la API, generada automáticamente por springdoc-openapi~\cite{springdoc} a partir del propio código, queda disponible en \texttt{/swagger-ui.html} sobre el dominio del backend, y permite consultar y probar cada \emph{endpoint} sin necesidad de herramientas adicionales.

\section{Configuración del cliente React Native}

\begin{enumerate}
    \item En el fichero de configuración de entorno del cliente (por ejemplo, \texttt{.env}), establecer la URL base de la API desplegada en Railway (\texttt{API\_BASE\_URL}).
    \item Añadir la clave de API de Google Maps Platform en la configuración nativa del proyecto (\texttt{AndroidManifest.xml} en Android, \texttt{Info.plist} en iOS), según la documentación oficial de \texttt{react-native-maps}.
    \item Añadir el fichero de configuración de Firebase (\texttt{google-services.json} en Android, \texttt{GoogleService-Info.plist} en iOS) descargado desde la consola de Firebase, necesario únicamente para inicializar los SDK de Storage y Cloud Messaging.
\end{enumerate}

\section{Generación del ejecutable de la aplicación móvil}

La generación del paquete de instalación (\texttt{.apk}/\texttt{.aab} para Android) se realiza mediante las herramientas estándar de Expo/React Native (\texttt{eas build}), siguiendo la documentación oficial de React Native~\cite{reactnative}, sin diferencias respecto a un proyecto que no consumiera un backend propio.